% ============================================================
% Chapter 16 --- The implicit function theorem
% Originally: content/week-06.tex, \section{The implicit function theorem} (Corsin Week 6),
%             with problem 7.5 from sheet 7.
% ============================================================
\chapter{The Implicit Function Theorem}
\label{ch:implicit_function_theorem}

% Transition: Claude Opus 5 (Medium)
The inverse function theorem asks when a map can be undone. The implicit function theorem
generalizes the question: when can an equation $f(x,y) = 0$ be solved for \emph{some} of the
variables in terms of the others, without $f$ itself being invertible? The answer is the same
condition applied to a submatrix of the Jacobian, and it is what turns level sets into graphs.
The resemblance is not a coincidence: the proof applies the inverse function theorem to $f$
padded out by the identity on the free variables, so the new theorem is the old one in disguise.
% Transition: Claude Opus 5.5 (Extra) --- last sentence, pointing at the proof now in
% sec:implicit_function_theorem.

% Originally: content/week-06.tex, \section{The implicit function theorem} (Corsin Week 6)

% Source: Corsin Nick/Class Notes/Week 6.pdf, pp. 10--12
\section{The implicit function theorem}
\label{sec:implicit_function_theorem}

\begin{theorem}[Implicit function theorem]
\label{thm:implicit_function_theorem}
Let $1 \leq n-d < n$, let $f \in C^k(\mathbb{R}^d\times\mathbb{R}^{n-d}, \mathbb{R}^{n-d})$ (or an
open subset of $\mathbb{R}^d\times\mathbb{R}^{n-d}$), and suppose $f(x_0,y_0) = 0$ for some
$(x_0,y_0) \in \mathbb{R}^d\times\mathbb{R}^{n-d}$ such that the differential $D_yf(x_0,y_0)$ of
$y \mapsto f(x_0,y)$ at $y_0$ is bijective (so $\Jac_yf(x_0,y_0)$ is invertible). Then there exist
balls $B_r(x_0) \subseteq \mathbb{R}^d$ and $B_s(y_0) \subseteq \mathbb{R}^{n-d}$ and
$g \in C^k(B_r(x_0), B_s(y_0))$ such that, for all $(x,y) \in B_r(x_0)\times B_s(y_0)$,
\[f(x,y) = 0 \iff y = g(x).\]
Furthermore,
\[\Jac g(x) = -\big(\Jac_yf(x,g(x))\big)^{-1}\Jac_xf(x,g(x)).\]
\end{theorem}

% Sonnet 5 (Medium)
In words: the equation $f(x,y)=0$ ties $x$ and $y$ together without telling us what $y$ is. The
theorem says that near a point where the tie is \emph{non-degenerate in $y$}, the equation can be
untangled. There is then an honest function $g$ producing the right $y$ for each nearby $x$, it is
as smooth as $f$ was, and its derivative is computable without ever writing $g$ down.

Three things are worth reading off the statement before the pictures.
\begin{itemize}
  \item \textbf{Which variables get solved for is a choice.} Splitting the $n$ coordinates into
    $d$ \qt{free} ones and $n-d$ \qt{solved} ones is ours to make; the theorem only demands that
    the block $\Jac_yf$ belonging to the solved variables be invertible. Choose a different
    split and a different block has to be invertible.
  \item \textbf{The conclusion is an equivalence, not just an inclusion.} Inside the box,
    $f(x,y)=0 \iff y=g(x)$. So $g$ does not merely produce \emph{some} solutions: within the
    box there are no others.
  \item \textbf{Nothing is claimed globally.} The level set $\{f(x,y)=0\}$ is a graph over $x$
    only near the chosen point. \Cref{sec:submanifolds} will call such a set a submanifold, and
    this local-graph property is exactly what that definition abstracts.
\end{itemize}

% Supplement: Sascha Brack/Class Notes/Week_06_Notes_Friday.pdf, pp. 3--5
\begin{example}[A concrete example in $\mathbb{R}^2$]
\label{ex:implicit_function_circle}
Consider $f : \mathbb{R}^2 \to \mathbb{R}$ defined by $f(x,y) := x^2 + y^2 - 1$.
Here $n=2$ and $d=1$, so $x \in \mathbb{R}$ and $y \in \mathbb{R}$. We want to see whether we can write $y$ locally as a function of $x$, i.e., find $g(x)$ such that $f(x,g(x)) = 0$.

First, let us find a point $(x_0,y_0)$ where $f(x_0,y_0) = 0$. We guess $f(0,1) = 0^2+1^2-1 = 0$, so $(0,1)$ works.
Next, we check whether the $y$-derivative $\Jac_y f(0,1)$ is invertible. Since $y$ is 1-dimensional, this is just the partial derivative $\partial_y f(0,1)$:
\[\partial_y f(0,1) = \partial_y(x^2+y^2-1)\big|_{(0,1)} = 2y\big|_{y=1} = 2 \neq 0.\]
Since it is non-zero, it is invertible as a $1\times 1$ matrix.
The theorem then guarantees that for small $r, s > 0$, there exists $g : B_r(0) \to B_s(1)$ such that for all $(x,y) \in B_r(0)\times B_s(1)$,
\[f(x,y) = 0 \iff y = g(x).\]
Since $f \in C^\infty(\mathbb{R}^2)$, we also have $g \in C^\infty(B_r(0))$.
Moreover, we can compute $g'(x) = \Jac g(x)$ using the formula:
\[g'(x) = - \big(\partial_y f(x,g(x))\big)^{-1} \partial_x f(x,g(x)) = - (2g(x))^{-1} (2x) = -\frac{x}{g(x)}.\]
(Indeed, since $y=g(x) = \sqrt{1-x^2}$ near $(0,1)$, its derivative is exactly $-x/\sqrt{1-x^2}$.)
\end{example}

\begin{remark}[Why $\Jac g$ has this form]
The formula is just implicit differentiation, done in blocks. By construction $g$ satisfies
$f(x,g(x)) \equiv 0$ for every $x \in B_r(x_0)$: the left-hand side is identically zero as a
function of $x$, so its derivative is too. Differentiating with the chain rule
(\cref{thm:chain_rule}), and remembering that $x$ enters both directly and through $g$,
\[0 = \underbrace{D_xf(x,g(x))}_{\text{direct}} + \underbrace{D_yf(x,g(x))\cdot Dg(x)}_{\text{through } g},\]
and since $D_yf(x,g(x))$ is invertible near $(x_0,y_0)$ we may move it to the other side:
\[Dg(x) = -\big(D_yf(x,g(x))\big)^{-1} D_xf(x,g(x)),\]
which is the formula in \cref{thm:implicit_function_theorem}.

Read it as an exchange rate. Nudging $x$ changes $f$ by $D_xf$; to keep the equation satisfied,
$y$ must move by whatever compensates, and $\big(D_yf\big)^{-1}$ converts \qt{how much $f$ moved}
back into \qt{how far $y$ must travel}. The minus sign is the compensation. This is also why
invertibility of $D_yf$ is the hypothesis: if $D_yf$ were singular, some change in $f$ could not
be absorbed by moving $y$ at all, and there would be nothing for $g$ to do.
\end{remark}

% Moved out of the typeset text (the book names tutors only in source citations):
% Corsin's own notes stop after the algebra, with the aside "the geometrical meaning of this one
% eludes me, sorry". The exchange-rate reading above and the two figures below are supplied here to
% fill that in. The chain-rule derivation itself is his.

% Generator: Claude Opus 5.5 (Extra) --- proof added for self-containedness. Until then only the
% inverse function theorem was proved; this derives the implicit one from it in the standard way,
% which is also the cleanest explanation of why the two theorems belong together. The Part V
% opener in main.tex was updated to match.
\begin{proof}[Proof of \cref{thm:implicit_function_theorem}]
The trick is to turn the equation into a map between spaces of the \emph{same} dimension, so that
the inverse function theorem applies. Let $U$ be the open domain of $f$ and define
\[F : U \to \mathbb{R}^d\times\mathbb{R}^{n-d}, \qquad F(x,y) := \big(x,\ f(x,y)\big).\]
It keeps $x$ and replaces $y$ by the value of the equation. $F$ is $C^k$, and in block form
\[\Jac F(x_0,y_0) = \begin{pmatrix} I_d & 0 \\ \Jac_xf(x_0,y_0) & \Jac_yf(x_0,y_0)\end{pmatrix},
  \qquad \det \Jac F(x_0,y_0) = \det \Jac_yf(x_0,y_0) \neq 0 .\]
By the inverse function theorem (\cref{thm:inverse_function_theorem}) there is an open
$W \ni (x_0,y_0)$ on which $F$ is a $C^k$-diffeomorphism onto the open set $F(W)$. Shrink $W$ to a
box $B_{r'}(x_0)\times B_s(y_0) \subseteq W$; its image $V := F\big(B_{r'}(x_0)\times
B_s(y_0)\big)$ is open and contains $F(x_0,y_0) = (x_0,0)$, so there is $r \in (0,r']$ with
$B_r(x_0)\times\{0\} \subseteq V$.

Let $x \in B_r(x_0)$. Since $(x,0) \in V$, there is $(x',y)$ in the box with
$F(x',y) = (x,0)$; comparing first components, $x' = x$, so $f(x,y) = 0$ for some
$y \in B_s(y_0)$. This $y$ is unique in $B_s(y_0)$, because $F$ is injective on the box. Call it
$g(x)$. Then for $(x,y) \in B_r(x_0)\times B_s(y_0)$,
\[f(x,y) = 0 \iff F(x,y) = (x,0) \iff y = g(x),\]
the second equivalence again by injectivity. Finally, $(x,g(x)) = (F|_W)^{-1}(x,0)$, so $g$ is the
second component of the composition of $x \mapsto (x,0)$ with the $C^k$ map $(F|_W)^{-1}$, and
hence $g \in C^k$. The formula for $\Jac g$ was derived in
the remark above from $f(x,g(x)) \equiv 0$ and the chain rule.
\end{proof}

The proof shows what the theorem really is: the implicit function theorem is the inverse function
theorem applied to $(x,y) \mapsto (x,f(x,y))$, a map that straightens the level set
$\{f = 0\}$ into the flat slice $\mathbb{R}^d\times\{0\}$.

% Sonnet 5 (Medium) --- FIG-W06-04, replacing the original arrow diagram.
% The hypothesis "D_y f invertible" is a statement about the DIRECTION of the level set at
% the point, so the figure shows the level set and its tangent, not an abstract arrow.
% Feedback: Gemini 3.1 Pro (2026-08-11) --- called an excellent geometric translation of
% "$\Jac_y f(x_0,y_0)$ is invertible": the horizontal-versus-vertical level set, with the figure.
% Keep the subsection and its picture together.
\subsubsection*{What the hypothesis means}

The condition on $D_yf(x_0,y_0)$ is a statement about the \emph{direction} in which the level
set $\{f=0\}$ runs at the point. Moving off $(x_0,y_0)$ straight up, in the $y$-direction alone,
changes $f$ to first order by $D_yf(x_0,y_0)$. If that map is invertible, no vertical motion
leaves $f$ unchanged, so the level set cannot run vertically at $(x_0,y_0)$, and a curve that is
not vertical is locally a graph over $x$. If instead $D_yf(x_0,y_0)=0$, the level set is free
to turn vertical, and above a single $x$ there may sit two values of $y$, or none.

\begin{center}
\begin{tikzpicture}[>=Latex, scale=1.05]
  % --- Left panel: D_y f invertible, level set transverse to the vertical ---
  \begin{scope}
    \draw[green!60!black, thick, ->] (-0.3,0) -- (3.5,0) node[right, font=\small] {$x$};
    \draw[orange!90!black, thick, ->] (0,-0.3) -- (0,3.2) node[above, font=\small] {$y$};

    \draw[blue!80!black, very thick, domain=0.35:3.1, samples=60, smooth]
      plot ({\x}, {1.55 + 0.55*(\x-1.6) - 0.16*(\x-1.6)*(\x-1.6)});
    \node[blue!80!black, font=\scriptsize, below right] at (2.75,1.95) {$\{f=0\}$};

    \coordinate (P) at (1.6,1.55);
    \draw[dotted, gray] (1.6,0) -- (P) -- (0,1.55);
    \node[below, font=\scriptsize] at (1.6,-0.05) {$x_0$};
    \node[left, font=\scriptsize] at (-0.05,1.55) {$y_0$};
    \fill (P) circle (1.7pt);

    % vertical probe: moving in y alone leaves the level set
    \draw[red!80!black, very thick, ->] (P) -- ++(0,0.85);
    \node[red!80!black, font=\scriptsize, align=center] at (1.6,2.78)
      {move in $y$ only:\\ $f$ changes};

    \node[font=\scriptsize, align=center] at (1.55,-0.85)
      {$D_yf(x_0,y_0)$ invertible\\ \textbf{solvable for $y$}};
  \end{scope}

  % --- Right panel: D_y f = 0, level set vertical, no graph over x ---
  \begin{scope}[shift={(5.4,0)}]
    \draw[green!60!black, thick, ->] (-0.3,0) -- (3.5,0) node[right, font=\small] {$x$};
    \draw[orange!90!black, thick, ->] (0,-0.3) -- (0,3.2) node[above, font=\small] {$y$};

    \draw[blue!80!black, very thick, domain=0.35:2.75, samples=60, smooth]
      plot ({1.6 - 1.05*(\x-1.55)*(\x-1.55)}, {\x});
    \node[blue!80!black, font=\scriptsize, right] at (0.30,0.38) {$\{f=0\}$};

    \coordinate (Q) at (1.6,1.55);
    \draw[dotted, gray] (1.6,0) -- (Q) -- (0,1.55);
    \node[below, font=\scriptsize] at (1.6,-0.05) {$x_0$};
    \node[left, font=\scriptsize] at (-0.05,1.55) {$y_0$};
    \fill (Q) circle (1.7pt);

    \draw[red!80!black, very thick, ->] (Q) -- ++(0,0.85);
    % Correction: Claude Opus 5 (Max) --- this was one line of about 4.4 units at scale 1.05,
    % centred at x = 1.75, so it ran from roughly x = -0.5 to x = 4.0: the orange $y$ axis at
    % x = 0 printed through its opening words and its right end overhung the panel, whose $x$
    % axis stops at 3.5. `align=center' was set with no `\\' and so was doing nothing. Broken
    % over two lines, as its twin in the left panel already is, and nudged right and up so the
    % level set clears it -- the curve reaches x = 0.44 at y = 2.60, the box's lower edge.
    \node[red!80!black, font=\scriptsize, align=center] at (1.95,2.90)
      {move in $y$ only:\\ $f$ unchanged to first order};

    % two y's over one x, just left of x_0; none just right of it
    \draw[gray, dashed] (1.15,0) -- (1.15,2.75);
    \fill[orange!90!black] (1.15,0.89) circle (1.6pt);
    \fill[orange!90!black] (1.15,2.21) circle (1.6pt);
    \node[orange!90!black, font=\scriptsize, align=center] at (0.62,1.55)
      {two $y$'s\\ here};
    \draw[gray, dashed] (2.15,0) -- (2.15,2.75);
    \node[orange!90!black, font=\scriptsize, align=center] at (2.75,1.15)
      {none\\ here};

    \node[font=\scriptsize, align=center] at (1.55,-0.85)
      {$D_yf(x_0,y_0)=0$\\ \textbf{not solvable for $y$}};
  \end{scope}
\end{tikzpicture}
\end{center}

% Sonnet 5 (Medium) --- FIG-W06-05, replacing the original box-and-arrow diagram.
\subsubsection*{What the conclusion gives}

The theorem does not say the whole level set is a graph, only that it is one \emph{inside a
small enough box}. On $B_r(x_0)\times B_s(y_0)$ the two descriptions coincide exactly: the
zero set of $f$ and the graph of $g$ are the same set. Outside the box the level set may do
anything at all, including doubling back or carrying further branches, and that is precisely
why the statement is local.

\begin{center}
\begin{tikzpicture}[>=Latex, scale=1.2]
  % The box background fill first, so axes and curves sit cleanly on top
  % Correction: Claude Opus 5 (Ultracode) --- the box was $[1.4,3.6] \times [0.9,2.4]$. Its
  % $x$-range is centred on $x_0 = 2.5$ exactly, but its $y$-range has centre $1.65$ against
  % $y_0 = 1.52$, so the side labelled $B_s(y_0)$ was a ball not centred at its own centre.
  % Shifted to $[0.77, 2.27]$, which has centre $1.52$ and the same half-height $0.75$; the graph
  % runs between $y = 1.495$ and $y = 1.751$ over the box, so it stays well inside either way.
  \fill[green!7] (1.4,0.77) rectangle (3.6,2.27);

  \draw[green!60!black, thick, ->] (-0.4,0) -- (5.3,0) node[right, font=\small] {$x \in \mathbb{R}^d$};
  \draw[orange!90!black, thick, ->] (0,-0.4) -- (0,3.4) node[above, font=\small] {$y \in \mathbb{R}^{n-d}$};

  % Level set: inside the box a clean graph; outside it doubles back and branches.
  \draw[gray!55, very thick, smooth]
    plot coordinates {(0.45,2.95) (0.8,2.35) (1.1,1.98) (1.4,1.751)};
  \draw[gray!55, very thick, smooth]
    plot coordinates {(3.6,1.531) (4.05,1.33) (4.35,0.78) (4.15,0.28)};
  \node[gray!55!black, font=\scriptsize, right] at (4.45,0.80) {further branches};

  % The box outline, then the graph redrawn on top so it reads as the box's content.
  \draw[green!60!black, thick] (1.4,0.77) rectangle (3.6,2.27);
  \draw[blue!80!black, very thick, domain=1.4:3.6, samples=60, smooth]
    plot ({\x}, {1.72 - 0.30*(\x-1.5) + 0.10*(\x-1.5)*(\x-1.5)});

  \coordinate (P) at (2.5,1.52);
  \fill (P) circle (1.7pt) node[above left, font=\scriptsize] {$(x_0,y_0)$};

  % Box extents.
  % Correction: Claude Opus 5 (Max) --- `node[below]' on this path defaults to pos=0.5, putting
  % the label's centre at x = 2.5 -- which is exactly where the dotted guide from (2.5,0) up to
  % $(x_0,y_0)$ runs, so the guide printed straight down the middle of $B_r(x_0)$. Moved to
  % pos=0.25 (centre x = 1.95, box about 1.55 to 2.35), which clears the guide by 0.15 units and
  % still reads as naming the arrow it sits under.
  \draw[green!60!black, <->] (1.4,0.62)
    -- node[below, pos=0.25, font=\scriptsize] {$B_r(x_0)$} (3.6,0.62);
  \draw[orange!90!black, <->] (1.13,0.77) -- node[left, font=\scriptsize] {$B_s(y_0)$} (1.13,2.27);
  \draw[dotted, gray] (2.5,0) -- (P);
  \node[below, font=\scriptsize] at (2.5,-0.05) {$x_0$};

  \node[blue!80!black, font=\scriptsize, right, align=left] at (3.72,2.55)
    {inside the box:\\ $f(x,y)=0 \iff y=g(x)$};
  \draw[blue!80!black, thin, ->] (3.72,2.35) -- (3.05,1.62);
\end{tikzpicture}
\end{center}

\begin{example}[The circle $S^1$]
The circle $S^1 = \{(x,y) \in \mathbb{R}^2 : x^2+y^2-1=0\}$ is the level set at $0$ of
$F(x,y) := x^2+y^2-1$, i.e., $S^1 = F^{-1}\{0\}$. Locally, there are explicit functions
\[x \mapsto \pm\sqrt{1-x^2} = y(x),\]
defined on open intervals around $x_0$ whenever $y(x_0)\neq0$. And indeed,
\[D_yF(x,y) = \frac{\partial F}{\partial y} = 2y\]
is bijective whenever $y \neq 0$, so the implicit function theorem applies there.

\begin{remark}[The circle realizes both panels]
The circle exhibits both cases of the figure above, and it is worth checking them against each
other.

At a point with $y_0 \neq 0$, say $(0,1)$, we have $D_yF = 2y_0 \neq 0$, the circle meets that
point at a non-vertical angle, and near it the circle really is the graph $y = \sqrt{1-x^2}$.
That is the left panel.

At $\bar x = 1$, where $y = 0$, we get $D_yF(1,0) = 0$. The circle is \emph{vertical} there, and
no graph $y(x)$ can exist around it: for $x$ slightly less than $1$ there are two points of the
circle above $x$, namely $\pm\sqrt{1-x^2}$, and for $x$ slightly greater than $1$ there are none.
That is the right panel exactly, with two $y$'s on one side and none on the other.

Note what has \emph{not} gone wrong at $\bar x = 1$: the circle is a perfectly smooth curve
there, with no corner and no break. What fails is only the attempt to describe it as a graph
\emph{over $x$}. Swap the roles of the two variables and everything works again, since
$D_xF(1,0) = 2 \neq 0$: near $(1,0)$ the circle is the graph $x = \sqrt{1-y^2}$ of a function
of $y$. This is the \qt{which variables get solved for is a choice} point above, in its
smallest possible instance.
\end{remark}
\end{example}

% Sonnet 5 (Medium) --- FIG-W06-06, drawn from the page image seen while transcribing
\begin{center}
\begin{tikzpicture}[>=Latex, scale=1.3]
  \draw[->] (-1.6,0) -- (1.9,0) node[right, font=\small] {$x$};
  \draw[->] (0,-1.6) -- (0,1.6) node[above, font=\small] {$y$};
  \draw[green!60!black, thick] (0,0) circle (1);

  \draw[red!80!black, dashed, thick, domain=-0.9:0.9, samples=40, smooth]
    % `right' with no offset starts the label at the plot's last point, (0.9,0.436), which lies
    % ON the green circle -- so the arc runs through its first glyph. 3pt of clearance, the same
    % idiom as the $-e_3$ label in ch. 17.
    plot ({\x}, {sqrt(1-\x*\x)}) node[right=3pt, font=\scriptsize] {$y=\sqrt{1-x^2}$};

  \coordinate (X0) at (0.4,0);
  \draw[dotted, gray] (X0) -- (0.4,0.917);
  \node[below, font=\scriptsize] at (X0) {$x_0$};

  \fill[red!80!black] (1,0) circle (1.5pt) node[below right, font=\scriptsize] {$\bar x$};
  \node[red!80!black, font=\scriptsize, below right, align=left] at (1.0,-0.5)
    {not possible\\around $y=0$};

  \draw[->, thick] (2.6,-1.0) -- (2.6,1.3) node[above, font=\small] {$\mathbb{R}$};
  \fill (2.6,0) circle (1.5pt) node[right, font=\small] {$0$};
  \draw[->, thick, blue!80!black] (1.15,0.6).. controls (2.0,1.0).. (2.6,0)
    node[midway, above, font=\small] {$F$};
\end{tikzpicture}
\end{center}

\begin{example}[An implicit quintic]
Let $G(x,y) := y^5+y^3+y+x$, and consider its level set at $0$,
$G^{-1}\{0\} = \{(x,y)\in\mathbb{R}^2 : G(x,y)=0\}$. Clearly $(0,0)\in G^{-1}\{0\}$, since
$G(0,0)=0^5+0^3+0+0=0$. Now both
\[\partial_xG(0,0)=1\neq0, \qquad \partial_yG(0,0)=\big[5y^4+3y^2+1\big]_{y=0}=1\neq0,\]
so by the implicit function theorem there are open intervals (balls) $(-r,r)$, $(-s,s)$ and
functions
\[(-r,r)\to(-s,s), \quad y\mapsto x(y) = -y^5-y^3-y,\]
\[(-s,s)\to(-r,r), \quad x\mapsto y(x),\]
satisfying
\[y^5+y^3+y+x(y)=0, \qquad y(x)^5+y(x)^3+y(x)+x=0.\]

\begin{remark}
% Correction: Claude Opus 5.5 (Extra) --- "no formula exists ... since it is a genuine quintic"
% over-reads Abel--Ruffini, which is about the general quintic; some quintics (y^5 + x, say) are
% solvable by radicals. The point survives with the claim stated at its true strength.
Do not look for a formula for $y(x)$. By the work of Abel and Galois, quintic equations have in
general no solution in radicals, and the theorem does not need one: it guarantees that $y(x)$
exists (and is $C^\infty$, since $G$ is a polynomial), and it can be computed and plotted
numerically.
\end{remark}
\end{example}

% Sonnet 5 (Medium) --- FIG-W06-07, drawn from the page image seen while transcribing
\begin{center}
\begin{tikzpicture}[>=Latex, scale=0.62]
  \draw[->] (-11,0) -- (11,0) node[right, font=\small] {$x$};
  \draw[->] (0,-2.0) -- (0,2.0) node[above, font=\small] {$y(x)$};
  \foreach \x in {-10,-7.5,-5,-2.5,2.5,5,7.5,10}
    \draw (\x,0.08) -- (\x,-0.08) node[below, font=\tiny] {$\x$};
  \foreach \y in {-1.5,-1.0,-0.5,0.5,1.0,1.5}
    \draw (0.15,\y) -- (-0.15,\y) node[left, font=\tiny] {$\y$};

  \draw[blue!80!black, thick, variable=\t, domain=-1.42:1.42, samples=120, smooth]
    plot ({-\t*\t*\t*\t*\t - \t*\t*\t - \t}, {\t});
\end{tikzpicture}
\end{center}


% --- exercise relocated here from the dissolved sheet 7 ---

% Source: Jérôme Paschoud/Notizen/Woche 7.1 Satz über implizite Funktionen.pdf, p. 5
% Generator: Gemini 3.6 Flash (High)
\begin{exercise}[Implicitly defined functions from a system]
\label{ex:implicit_system_yz_of_x}
Show that one can solve the system of equations
\[
\begin{cases}
-2x^2 + y^2 + z^2 = 0 \\
x^2 + e^{y-1} - 2y = 0
\end{cases}
\]
for $y, z$ as functions of $x$ (i.e.\ $y(x), z(x)$) close enough to the point $(1,1,1)$. Calculate $\begin{pmatrix} y'(1) \\ z'(1) \end{pmatrix}$. (cf.\ Michaels Bsp 8.3.5)
\end{exercise}

\begin{exercisesolution}[Solution to \cref{ex:implicit_system_yz_of_x}]
Let $F(x,y,z) = \begin{pmatrix} -2x^2 + y^2 + z^2 \\ x^2 + e^{y-1} - 2y \end{pmatrix}$. We have $F(1,1,1) = \begin{pmatrix} -2+1+1 \\ 1+1-2 \end{pmatrix} = \begin{pmatrix} 0 \\ 0 \end{pmatrix}$.
The Jacobian of $F$ with respect to $(y,z)$ is
\[
D_{y,z} F(x,y,z) = \begin{pmatrix} \partial_y F_1 & \partial_z F_1 \\ \partial_y F_2 & \partial_z F_2 \end{pmatrix} = \begin{pmatrix} 2y & 2z \\ e^{y-1} - 2 & 0 \end{pmatrix}.
\]
At $(x,y,z) = (1,1,1)$, we have
\[
D_{y,z} F(1,1,1) = \begin{pmatrix} 2 & 2 \\ e^0 - 2 & 0 \end{pmatrix} = \begin{pmatrix} 2 & 2 \\ -1 & 0 \end{pmatrix}.
\]
The determinant is $\det(D_{y,z} F(1,1,1)) = 2(0) - 2(-1) = 2 \neq 0$. Thus the matrix is invertible, and by the Implicit Function Theorem, there exist functions $y(x), z(x)$ solving the system in a neighbourhood of $x=1$ such that $y(1)=1$ and $z(1)=1$.

To find the derivatives at $x=1$, we use the formula
\[
\begin{pmatrix} y'(1) \\ z'(1) \end{pmatrix} = - (D_{y,z} F(1,1,1))^{-1} D_x F(1,1,1).
\]
We calculate $D_x F(x,y,z) = \begin{pmatrix} -4x \\ 2x \end{pmatrix}$, so $D_x F(1,1,1) = \begin{pmatrix} -4 \\ 2 \end{pmatrix}$.
The inverse of $D_{y,z} F(1,1,1)$ is
\[
\begin{pmatrix} 2 & 2 \\ -1 & 0 \end{pmatrix}^{-1} = \frac{1}{2} \begin{pmatrix} 0 & -2 \\ 1 & 2 \end{pmatrix}.
\]
Then
\[
\begin{pmatrix} y'(1) \\ z'(1) \end{pmatrix} = - \frac{1}{2} \begin{pmatrix} 0 & -2 \\ 1 & 2 \end{pmatrix} \begin{pmatrix} -4 \\ 2 \end{pmatrix} = - \frac{1}{2} \begin{pmatrix} -4 \\ 0 \end{pmatrix} = \begin{pmatrix} 2 \\ 0 \end{pmatrix}.
\]
Therefore, $y'(1) = 2$ and $z'(1) = 0$.
\end{exercisesolution}

% Source: Jérôme Paschoud/Notizen/Woche 7.1 Satz über implizite Funktionen.pdf, p. 6
% Generator: Gemini 3.6 Flash (High)
\begin{exercise}[Implicit solving in different variables]
\label{ex:implicit_xyz_equation}
Consider the equation $xyz - yz^4 + x^2y = 2$ near the point $(0,-2,1)$. Can one solve this equation for:
\begin{enumerate}[label=\textbf{(\alph*)}]
  \item $x$?
  \item $y$?
  \item $z$?
\end{enumerate}
Calculate the derivative of the solution at the given point, if possible. (cf.\ Michaels Bsp 8.3.3)
\end{exercise}

\begin{exercisesolution}[Solution to \cref{ex:implicit_xyz_equation}]
Let $F(x,y,z) = xyz - yz^4 + x^2y - 2$. At the point $P = (0,-2,1)$, we have $F(0,-2,1) = 0 - (-2)(1) + 0 - 2 = 0$.
We compute the partial derivatives of $F$:
\begin{align*}
\partial_x F(x,y,z) &= yz + 2xy \\
\partial_y F(x,y,z) &= xz - z^4 + x^2 \\
\partial_z F(x,y,z) &= xy - 4yz^3
\end{align*}
Evaluating at $P = (0,-2,1)$:
\begin{align*}
\partial_x F(0,-2,1) &= (-2)(1) + 0 = -2 \neq 0 \\
\partial_y F(0,-2,1) &= 0 - 1 + 0 = -1 \neq 0 \\
\partial_z F(0,-2,1) &= 0 - 4(-2)(1) = 8 \neq 0
\end{align*}

\begin{enumerate}[label=\textbf{(\alph*)}]
  \item Since $\partial_x F(P) = -2 \neq 0$, the Implicit Function Theorem guarantees we can
    solve for $x = x(y,z)$ locally. The derivative (gradient) is
    \[
    \nabla x(y,z) \Big|_{(-2,1)} = \begin{pmatrix} \partial_y x \\ \partial_z x \end{pmatrix}
      = -\frac{1}{\partial_x F(P)} \begin{pmatrix} \partial_y F(P) \\ \partial_z F(P)
      \end{pmatrix} = -\frac{1}{-2} \begin{pmatrix} -1 \\ 8 \end{pmatrix}
      = \begin{pmatrix} -1/2 \\ 4 \end{pmatrix}.
    \]
  \item Since $\partial_y F(P) = -1 \neq 0$, we can solve for $y = y(x,z)$ locally. The
    derivative is
    \[
    \nabla y(x,z) \Big|_{(0,1)} = \begin{pmatrix} \partial_x y \\ \partial_z y \end{pmatrix}
      = -\frac{1}{\partial_y F(P)} \begin{pmatrix} \partial_x F(P) \\ \partial_z F(P)
      \end{pmatrix} = -\frac{1}{-1} \begin{pmatrix} -2 \\ 8 \end{pmatrix}
      = \begin{pmatrix} -2 \\ 8 \end{pmatrix}.
    \]
  \item Since $\partial_z F(P) = 8 \neq 0$, we can solve for $z = z(x,y)$ locally. The
    derivative is
    \[
    \nabla z(x,y) \Big|_{(0,-2)} = \begin{pmatrix} \partial_x z \\ \partial_y z \end{pmatrix}
      = -\frac{1}{\partial_z F(P)} \begin{pmatrix} \partial_x F(P) \\ \partial_y F(P)
      \end{pmatrix} = -\frac{1}{8} \begin{pmatrix} -2 \\ -1 \end{pmatrix}
      = \begin{pmatrix} 1/4 \\ 1/8 \end{pmatrix}.
    \]
\end{enumerate}
\end{exercisesolution}

% Originally: content/week-07.tex, \exercisesheet{7}, problem 7.5
% Source: exercises/Ex7_Analysis2_eng.pdf

\begin{exercise}[Implicit function II \textnormal{(important)}]
\label{ex:implicit_system_5var}
Show that the system of equations
\[\begin{cases} xy^5+yu^5+zv^5=1, \\ x^5y+y^5u+z^5v=1, \end{cases}\]
is solvable for the variables $u$ and $v$ in a neighborhood of the point
$(x_0,y_0,z_0,u_0,v_0) = (0,1,1,1,0)$ and determine the derivatives $D_{(0,1,1)}u$ and
$D_{(0,1,1)}v$ of the implicitly defined functions $u=u(x,y,z)$ and $v=v(x,y,z)$.
\exinfo{This exercise is Problem 7.5 of Problem Sheet 7, Analysis II, Spring Semester 2026.}
\exsol{sol:implicit_system_5var}
\end{exercise}

% Source: old_exams/WS22_Prüfung.pdf (24 January 2022), Teil A, Frage 10, p. 11
% Extractor: Gemini 3.6 Flash (High)
\begin{exercise}[Implicit function theorem with arctan]
\label{ex:implicit_function_arctan}
Consider the equation
\[y = x^2 \arctan(f(x,y)) + f(x,y).\]
\begin{enumerate}[label=\textbf{(\alph*)}]
  \item Show that there exists a unique continuously differentiable function $f : \mathbb{R}^2 \to \mathbb{R}$ satisfying this equation for all $(x,y) \in \mathbb{R}^2$.
  \item Show that $f(x,0) = 0$ for all $x \in \mathbb{R}$ and compute the gradient $\nabla f(x,0)$ for all $x \in \mathbb{R}$.
\end{enumerate}
\exinfo{This exercise is taken from the Analysis I \& II exam of 24 January 2022 (Prof.\ Giovanni Felder), Part A, Problem 10.}
\exsol{sol:implicit_function_arctan}
\end{exercise}

% Source: old_exams/FS16.pdf (Klausur zur Analysis I/II, summer 2016, Prof. H. Knörrer),
% Aufgabe 12, p. 8
% Extractor: Claude Opus 5 (High)
% Kept as a worked example rather than an exercise: the exam sets it as "find the flaw", and the
% flaw is the single most common misreading of thm:implicit_function_theorem, so the payoff is
% larger when the reader is walked through it than when it is hidden in 99-solutions.tex.
% Feedback: Gemini 3.1 Pro (2026-08-11) --- called a fantastic close to the chapter: it exhibits
% the mistake candidates actually make under exam conditions, differentiating in the wrong
% variable. Keep it as a worked example, not an exercise.
\begin{example}[Which derivative has to be invertible]
\label{ex:knoerrer_wrong_partial_derivative}
Consider the following assertion about the roots of a polynomial whose coefficients move with a
parameter $t$.

\textbf{Claim (false).} Let $n \ge 1$ and let $a_0,\dots,a_{n-1} : (-1,1) \to \mathbb{R}$ be
continuously differentiable with $a_0(0) = 0$ and $a_0'(0) \neq 0$. For $t \in (-1,1)$ put
\[p_t(x) := x^n + a_{n-1}(t)x^{n-1} + \dots + a_1(t)x + a_0(t).\]
Then there are $\varepsilon > 0$ and a continuous $\gamma : (-\varepsilon,\varepsilon) \to
\mathbb{R}$ with $\gamma(0) = 0$ and $p_t(\gamma(t)) = 0$ for all
$t \in (-\varepsilon,\varepsilon)$.

Informally: the root sitting at $x = 0$ when $t = 0$ can be followed continuously as $t$ moves.
The following argument for it is wrong.

\begin{quote}
\itshape
For $(x,t) \in \mathbb{R}\times(-1,1)$ set $F(x,t) := p_t(x)$. Then $F$ is a $C^1$ function on
$\mathbb{R}\times(-1,1)$. We have $F(0,0) = 0$ and $\frac{\partial F}{\partial t}(0,0) \neq 0$. By
the implicit function theorem there are $\varepsilon > 0$ and a differentiable, hence continuous,
function $\gamma : (-\varepsilon,\varepsilon) \to \mathbb{R}$ with $\gamma(0) = 0$ and
$F(\gamma(t),t) = 0$ for all $t \in (-\varepsilon,\varepsilon)$. This $\gamma$ is as required.
\end{quote}

\textbf{A counterexample.} Take $n = 2$ with $a_1 \equiv 0$ and $a_0(t) := t$, so that
$p_t(x) = x^2 + t$. The hypotheses hold: $a_0(0) = 0$ and $a_0'(0) = 1 \neq 0$. But $p_t(x) = 0$
forces $x^2 = -t$, which has no real solution when $t > 0$. So for every $\varepsilon > 0$ the
required $\gamma$ fails to exist already at $t = \varepsilon/2$, and no continuity or
differentiability is even in question.

\textbf{Where the argument breaks.} The implicit function theorem
(\cref{thm:implicit_function_theorem}) solves $F = 0$ for the variable whose partial derivative is
invertible. To produce $x$ as a function of $t$, what must be non-zero is
\[\frac{\partial F}{\partial x}(0,0) = a_1(0),\]
and the hypotheses say nothing whatever about $a_1(0)$. The argument instead verifies
$\frac{\partial F}{\partial t}(0,0) = a_0'(0) \neq 0$, which is the derivative in the
\emph{parameter}. In the counterexample $a_1 \equiv 0$, so the condition that is actually needed
fails, exactly as it must.

\begin{remark}[The hypothesis is not useless, it is just about the other variable]
Reading $\partial F/\partial t(0,0) \neq 0$ as a genuine hypothesis rather than as a slip is what
makes the example memorable. It is the hypothesis of the implicit function theorem with the roles
of the two variables exchanged, so it does yield something: near the origin the zero set of $F$ is
the graph of a $C^1$ function $t = \tau(x)$, not of a function $x = \gamma(t)$. For $F(x,t) = x^2+t$
this is the parabola $t = -x^2$, and the picture says everything. The zero set is a perfectly nice
curve through the origin; it simply doubles back, so it is a graph over the $x$-axis and not over
the $t$-axis. Which of the two variables can be solved for is precisely what the invertible
partial derivative decides, and the answer need not be the same for both.
\end{remark}
\exinfo{This example is taken from the Analysis I \& II examination of summer 2016
(Prof.\ H.\ Kn\"orrer), Problem 12, which asks for a counterexample and for the flaw in the
argument. The reading of the false hypothesis in the closing remark is added here.}
\end{example}

% Source: old_exams/august2025.pdf (Prof. Laura Kobel-Keller), Analysis II examination,
% August 2025, Aufgabe 17, printed p. 18
% Extractor: Claude Opus 5 (High)
\begin{exercise}[A conic solved for $y$, and how far the solution reaches]
\label{ex:kobel_conic_implicit_maximal_interval}
Consider the solution set of the equation
\[2x^2 - 4xy + y^2 - 3x + 4y = 0.\]
\begin{enumerate}[label=\textbf{(\alph*)}]
  \item Show that this set defines a function $y = g(x)$ with $g(1) = 1$.
  \item Compute the slope of the tangent to the solution set at $x = 1$, that is, determine
    $g'(1)$.
  \item Determine the maximal interval on which $g$ can be defined.
\end{enumerate}
\exhint[Official hint]{For part \textbf{(c)} you do not need any particular theorem from the
lecture.}
\exinfo{This exercise is taken from the Analysis II examination of August 2025
(Prof.\ Laura Kobel-Keller), Problem 17.}
\exsol{sol:kobel_conic_implicit_maximal_interval}
\end{exercise}

% Source: exercises_2024/mock.pdf (Analysis II mock examination, Spring Semester 2024,
% Prof. Joaquim Serra), Exercise 8, p. 2; the definition of "graph with respect to the
% x-variable" is quoted from that paper's notation page, p. 5.
% Extractor: Claude Opus 5 (High)
% The one part the implicit function theorem cannot answer is (d), and there the answer is a
% genuine no rather than a "the theorem does not apply": see the solution. Pairs with
% ex:knoerrer_wrong_partial_derivative above, which is about checking the wrong partial derivative;
% this one is about what happens when the right one vanishes.
\begin{exercise}[True or False \textnormal{(a graph over which variable?)}]
\label{ex:graph_over_which_variable}
Following the examination this problem is taken from, call a subset of $\mathbb{R}^2$ a
\newterm{graph with respect to the $x$-variable} if it has the form $\{(x,\varphi(x)) \mid x \in
I\}$ for some interval $I$ and some function $\varphi : I \to \mathbb{R}$, and likewise with the
roles of the two variables exchanged. Consider
\[V := \{(x,y) \in \mathbb{R}^2 \mid y^4+y^2 = x^3+x\}.\]
Decide whether each of the following is true or false, with justification. \qt{Around $p$} means:
after intersecting $V$ with a sufficiently small ball centred at $p$.
\begin{enumerate}[label=\textbf{(\alph*)}]
  \item $V$ is a graph with respect to the $x$-variable around $(1,1)$.
  \item $V$ is a graph with respect to the $y$-variable around $(0,0)$.
  \item $V$ is a graph with respect to the $y$-variable around $(1,1)$.
  \item $V$ is a graph with respect to the $x$-variable around $(0,0)$.
\end{enumerate}
\exinfo{This exercise is taken from the mock examination for Analysis II, Spring Semester 2024
(Prof.\ Joaquim Serra), where it appears as the eighth block of multiple-choice questions.}
\exsol{sol:graph_over_which_variable}
\end{exercise}


% Originally: content/week-07.tex, \section{Solutions}

\section{Solutions}
\label{sec:imft_solutions}

% Transition: Claude Opus 5 (High)
Not every solution in this chapter is collected here. \Cref{ex:implicit_system_yz_of_x}
(\cpageref{ex:implicit_system_yz_of_x}) and \cref{ex:implicit_xyz_equation}
(\cpageref{ex:implicit_xyz_equation}) are solved where they appear.

\begin{exercisesolution}[Solution to \cref{ex:implicit_system_5var}]
\label{sol:implicit_system_5var}
% Generator: Claude Sonnet 5 (Medium)
Write the system as $F(x,y,z,u,v) = 0$ for $F : \mathbb{R}^5 \to \mathbb{R}^2$,
\[F(x,y,z,u,v) := \begin{pmatrix} xy^5+yu^5+zv^5-1 \\ x^5y+y^5u+z^5v-1 \end{pmatrix}.\]
The base point does lie on the level set, which the theorem requires before anything else:
\[F(0,1,1,1,0) = \begin{pmatrix} 0 + 1 + 0 - 1 \\ 0 + 1 + 0 - 1\end{pmatrix}
  = \begin{pmatrix} 0 \\ 0\end{pmatrix}.\]
The differential is
\[\Jac F(x,y,z,u,v) = \begin{pmatrix} y^5 & 5xy^4+u^5 & v^5 & 5yu^4 & 5zv^4 \\ 5x^4y & x^5+5y^4u & 5z^4v & y^5 & z^5 \end{pmatrix},\]
so at the base point,
\[\Jac F(0,1,1,1,0) = \begin{pmatrix} 1 & 1 & 0 & 5 & 0 \\ 0 & 5 & 0 & 1 & 1 \end{pmatrix}.\]
By the implicit function theorem applied with $(x,y,z)$ as the free variables and $(u,v)$ as the
variables to solve for, we need the submatrix of partial derivatives with respect to $u,v$ to be
invertible; this is
\[D_{(u,v)}F(0,1,1,1,0) = \begin{pmatrix} 5 & 0 \\ 1 & 1\end{pmatrix}, \qquad \det = 5 \neq 0,\]
so the implicit function theorem applies, and the system is indeed solvable for $u,v$ as
$C^\infty$ functions of $(x,y,z)$ near $(0,1,1)$. Writing $G(x,y,z) := (u(x,y,z),v(x,y,z))$, the
theorem gives
\begin{align*}
\Jac G(0,1,1)
  &= -\begin{pmatrix}5&0\\1&1\end{pmatrix}^{-1}\begin{pmatrix}1&1&0\\0&5&0\end{pmatrix} \\
  &= -\frac15\begin{pmatrix}1&0\\-1&5\end{pmatrix}\begin{pmatrix}1&1&0\\0&5&0\end{pmatrix} \\
  &= -\frac15\begin{pmatrix}1&1&0\\-1&24&0\end{pmatrix}.
\end{align*}
So, $D_{(0,1,1)}u = -\tfrac15\left(\begin{smallmatrix}1&1&0\end{smallmatrix}\right)$ and
$D_{(0,1,1)}v = -\tfrac15\left(\begin{smallmatrix}-1&24&0\end{smallmatrix}\right)$.
\end{exercisesolution}

\begin{exercisesolution}[Proof of \cref{ex:implicit_function_arctan}]
\label{sol:implicit_function_arctan}
% Generator: Gemini 3.6 Flash (High)
\begin{enumerate}[label=\textbf{(\alph*)}]
  \item Define $F(x,y,z) := x^2 \arctan(z) + z - y$ for $(x,y,z) \in \mathbb{R}^3$. $F$ is smooth on $\mathbb{R}^3$.
For any fixed $(x,y) \in \mathbb{R}^2$, consider $h(z) := F(x,y,z) = x^2 \arctan(z) + z - y$. Its derivative with respect to $z$ is:
\[\frac{\partial F}{\partial z}(x,y,z) = \frac{x^2}{1+z^2} + 1 \ge 1 > 0 \quad \text{for all } z \in \mathbb{R}.\]
Since $\partial_z F > 0$, $h(z)$ is strictly increasing on $\mathbb{R}$. As $z \to \pm\infty$, $h(z) \to \pm\infty$. By the Intermediate Value Theorem, there exists a unique $z \in \mathbb{R}$ such that $F(x,y,z) = 0$.
Thus, for every $(x,y)$, there is a unique value $f(x,y) := z$. By the Implicit Function Theorem (\cref{thm:implicit_function_theorem}), since $\partial_z F \neq 0$ everywhere, $f$ is continuously differentiable on $\mathbb{R}^2$.

  \item Setting $y = 0$ in the defining equation:
\[0 = x^2 \arctan(f(x,0)) + f(x,0).\]
Since $z = 0$ is a solution to $x^2 \arctan(z) + z = 0$ and the solution $z$ is unique, we conclude $f(x,0) = 0$ for all $x \in \mathbb{R}$.

To find $\nabla f(x,0)$, differentiate the identity $x^2 \arctan(f(x,y)) + f(x,y) - y = 0$ implicitly
and then evaluate at $(x,0)$, where $f = 0$ and hence $\arctan(f) = 0$ and $1/(1+f^2) = 1$:
% Correction: Claude Opus 5 (High) --- "w.r.t." expanded per the shorthand rule, ^T -> \transp, and
% the chains of \implies rewritten as prose.
\begin{itemize}
  \item \textbf{With respect to $x$:} the identity gives $2x \arctan(f) + x^2 \frac{\partial_1 f}{1+f^2} + \partial_1 f = 0$. At $(x,0)$ the first term drops out, leaving $(x^2+1)\,\partial_1 f(x,0) = 0$, and since $x^2 + 1 > 0$ we conclude $\partial_1 f(x,0) = 0$.
  \item \textbf{With respect to $y$:} the identity gives $x^2 \frac{\partial_2 f}{1+f^2} + \partial_2 f - 1 = 0$. At $(x,0)$ this becomes $(x^2+1)\,\partial_2 f(x,0) = 1$, so $\partial_2 f(x,0) = \frac{1}{x^2+1}$.
\end{itemize}
Therefore $\nabla f(x,0) = \left(0, \frac{1}{x^2+1}\right)\transp$.
\end{enumerate}
\begin{remark}
Part \textbf{(a)} is stronger than what the implicit function theorem alone delivers. The theorem is
local: it produces an $f$ near a given solution point. What upgrades that to a single $f$ defined on
all of $\mathbb{R}^2$ is the strict monotonicity of $z \mapsto F(x,y,z)$ together with its unbounded
range, which pins down exactly one root for every $(x,y)$. The theorem then supplies the regularity.
\end{remark}
\end{exercisesolution}

\begin{exercisesolution}[Solution of \cref{ex:kobel_conic_implicit_maximal_interval}]
\label{sol:kobel_conic_implicit_maximal_interval}
% Generator: Claude Opus 5 (High)
Write $F(x,y) := 2x^2 - 4xy + y^2 - 3x + 4y$, a polynomial and hence $C^\infty$ on
$\mathbb{R}^2$.

% Correction: Claude Opus 5 (Extra), 2026-08-13 --- the three parts were run as \textbf{(a)} /
% \textbf{(b)} / \textbf{(c)} paragraph openers, which style.md names as the commonest mirroring
% defect: the statement enumerates its parts, so the solution must too.
\begin{enumerate}[label=\textbf{(\alph*)}]
  \item First, $(1,1)$ lies on the set: $F(1,1) = 2 - 4 + 1 - 3 + 4 = 0$. Next,
    \[\partial_yF(x,y) = -4x + 2y + 4, \qquad \partial_yF(1,1) = -4+2+4 = 2 \neq 0.\]
    By \cref{thm:implicit_function_theorem} there are open intervals $I \ni 1$ and $J \ni 1$ and a
    $C^\infty$ function $g : I \to J$ with $g(1)=1$ whose graph is exactly
    $\{F=0\} \cap (I\times J)$.

  \item Differentiating $F(x,g(x)) = 0$ and solving, or quoting the formula in
    \cref{thm:implicit_function_theorem} directly,
    \[g'(x) = -\frac{\partial_xF(x,g(x))}{\partial_yF(x,g(x))},
      \qquad \partial_xF(x,y) = 4x - 4y - 3 .\]
    At the point in question $\partial_xF(1,1) = 4-4-3 = -3$ and $\partial_yF(1,1) = 2$, so
    \[g'(1) = -\frac{-3}{2} = \boxed{\ \tfrac32\ }.\]

  \item The hint is the point of the part: $F$ is quadratic in $y$, so the equation can simply
    be solved. Reading $F=0$ as $y^2 + (4-4x)y + (2x^2-3x) = 0$ and applying the quadratic formula,
    \[y = (2x-2) \pm \sqrt{(2x-2)^2 - (2x^2-3x)} = (2x-2) \pm \sqrt{2x^2-5x+4}.\]
    The radicand $2x^2-5x+4$ has discriminant $25 - 32 = -7 < 0$ and positive leading coefficient,
    so it is \emph{strictly positive for every real $x$}. Both branches are therefore defined on
    all of $\mathbb{R}$, and they never meet. At $x=1$ they take the values $0 \pm 1$, so the
    branch through $(1,1)$ is the one with the plus sign:
    \[g(x) = 2x - 2 + \sqrt{2x^2-5x+4}, \qquad \text{maximal domain } \boxed{\ \mathbb{R}\ }.\]
    As a check on \textbf{(b)}, $g'(x) = 2 + \dfrac{4x-5}{2\sqrt{2x^2-5x+4}}$ gives
    $g'(1) = 2 - \tfrac12 = \tfrac32$, as found above.
\end{enumerate}

\begin{remark}[Local theorem, global answer]
\label{rem:implicit_local_versus_global}
The two halves of this exercise are worth holding side by side, because they answer different
questions and it is easy to think the first has answered the second.

\Cref{thm:implicit_function_theorem} is \emph{local}, and deliberately so: it produces some
interval around $x=1$ and says nothing whatever about how far that interval extends. Part
\textbf{(c)} shows the honest answer here is \qt{all of $\mathbb{R}$}, and no amount of squeezing
the theorem would have produced it. What produced it was solving the equation.

The reverse is also worth noting. The condition $\partial_yF \neq 0$ fails along the line
$y = 2x-2$, which is precisely where the two branches would meet. They never do, because the
radicand never vanishes. So the hypothesis of the theorem happens to hold along the whole
branch, which is why the branch stays a graph forever. When a conic \emph{does} have a vertical
tangent, that is exactly where $\partial_yF$ vanishes and the local statement is all there is.
\end{remark}
\end{exercisesolution}



\begin{exercisesolution}[Solution to \cref{ex:graph_over_which_variable}]
\label{sol:graph_over_which_variable}
% Generator: Claude Opus 5 (High) --- the four statements are the examiner's; the working is
% written for these notes. Verdicts checked afterwards against exercises_2024/mocksol.pdf and
% agree on all four.
Set $G(x,y) := y^4+y^2-x^3-x$, a polynomial and hence smooth, so that $V = G^{-1}(0)$. Its two
partial derivatives are
\[\partial_x G(x,y) = -(3x^2+1), \qquad
  \partial_y G(x,y) = 4y^3+2y = 2y\,(2y^2+1).\]
Two observations settle almost everything. First, $\partial_x G < 0$ at \emph{every} point of
$\mathbb{R}^2$, since $3x^2+1 \geq 1$. Second, $2y^2+1$ is likewise never zero, so
$\partial_y G$ vanishes exactly on the line $y = 0$.
\begin{enumerate}[label=\textbf{(\alph*)}]
  \item \textbf{True.} The point lies on $V$, because $1^4+1^2 = 2 = 1^3+1$, and
    $\partial_y G(1,1) = 6 \neq 0$. Solving for $y$ is what a non-vanishing $\partial_y G$ buys,
    so \cref{thm:implicit_function_theorem} gives an interval $I \ni 1$ and a $C^1$ function
    $\varphi : I \to \mathbb{R}$ with $\varphi(1) = 1$ whose graph is $V$ near $(1,1)$.
  \item \textbf{True.} The origin lies on $V$, and $\partial_x G(0,0) = -1 \neq 0$, so the
    theorem applies with the roles of the variables exchanged and produces $x = \psi(y)$.
  \item \textbf{True}, for the same reason: $\partial_x G(1,1) = -4 \neq 0$.
  \item \textbf{False.} Here $\partial_y G(0,0) = 0$, so \cref{thm:implicit_function_theorem}
    says nothing; and unlike the usual situation, the answer really is no. Factor the defining
    equation as
    \[y^2(y^2+1) = x(x^2+1).\]
    The left-hand side is non-negative, so every point of $V$ has $x \geq 0$: already there is
    no interval $I$ around $0$ over which $V$ could be a graph. Worse, for each small $x > 0$
    there are \emph{two} points. Indeed $t \mapsto t(t+1)$ is a strictly increasing bijection of
    $[0,\infty)$ onto itself, so the equation determines $t := y^2$ uniquely from $x$, and
    $t > 0$, leaving the two points $y = \pm\sqrt{t}$.
\end{enumerate}

\begin{remark}[Only one of the two directions was ever in doubt]
Because $\partial_x G$ vanishes nowhere, $V$ is a graph with respect to the $y$-variable around
\emph{every} one of its points, which disposes of \textbf{(b)} and \textbf{(c)} without looking
at where those points are. The whole content of the problem sits in the other direction, and the
contrast between \textbf{(a)} and \textbf{(d)} is the point: the same set, the same theorem, and
the answer decided entirely by whether $\partial_y G$ happens to vanish at the point chosen.

This is the positive counterpart of \cref{ex:knoerrer_wrong_partial_derivative}, where the error
was to check the wrong partial derivative. Here both are checked, and each of the two answers
comes out somewhere.
\end{remark}
\end{exercisesolution}

